% GENERATED FILE. Edit canonical lessons, then run npm run generate:books.
% canonical-source-hash: fnv1a64:54dc413afcd56339
% canonical-lessons: HI-C349-ugalna, HI-C349-inkarkarna, HI-C349-pustikarna, HI-C349-svikarkarna, HI-C349-galatfahmi

\chapter{Denial and Confession}
\label{ch:hi-denial-and-confession}
\hlchaptermodality{349}

\begin{chapteropening}
\textbf{By the end of this chapter:} \emph{I can say who denied something and who owned up.}

Complete the last lesson of chapter 349: I can say who denied something and who owned up.
\end{chapteropening}

\section[\texorpdfstring{\hi{उगलना} (ugalnā)}{उगलना (ugalnā)}]{\hi{उगलना} (ugalnā) — to spit out; to blurt out (a secret, the truth)}

\label{lesson:HI-C349-ugalna}



\begin{quote}
Before the new one: say the Hindi for conclusive, then the Hindi for an accusation.
\end{quote}

\subsection*{You'll want to know: \hi{उगलना}}
\textbf{\hi{उगलना}} — \emph{ugalnā} — \textquotedblleft{}to spit out; to blurt out (a secret, the truth)\textquotedblright{}.

\begin{quote}
\hi{आख़िर} \hi{चोर} \hi{ने} \hi{सच} \hi{उगल} \hi{दिया।}
\end{quote}

\begin{quote}
\emph{ākhir cor ne sac ugal diyā.}
\end{quote}

\begin{quote}
In the end the thief came out with the truth.
\end{quote}

Past, with \hi{ने} (he or she …): \hi{उसने} \hi{सच} \hi{उगल} \hi{दिया}, \emph{usne sac ugal diyā}.

\begin{cousinweb}
\textbf{\hi{उगलना}} (\emph{ugalnā}) is a tadbhava, inherited from Sanskrit \emph{udgirati}, \textquotedblleft{}spits out\textquotedblright{} (CDIAL 1960).
\end{cousinweb}

\begin{grammarlens}[title={What you've built}]
One more word for giving reasons and opinions.
\end{grammarlens}

\subsection*{Your turn}
\begin{itemize}
\raggedright
  \item \emph{Say it:} \emph{ugalnā}
  \item \emph{Say it:} \emph{ugalnā}, once more
  \item \emph{Say it:} \emph{ākhir cor ne sac ugal diyā.} — In the end the thief came out with the truth.
  \item \emph{Recall:} say the Hindi for conclusive, then the Hindi for an accusation, then say \emph{ugalnā} again
\end{itemize}

\begin{tcolorbox}[breakable,colback=teal!4,colframe=teal!35!black,title={Before you move on}]
What does \hi{उगलना} mean? (\textquotedblleft{}To spit out; to blurt out (a secret, the truth)\textquotedblright{}.) Say it once more.
\end{tcolorbox}
\section[\texorpdfstring{\hi{इनकार} \hi{करना} (inkār karnā)}{इनकार करना (inkār karnā)}]{\hi{इनकार} \hi{करना} (inkār karnā) — to deny, to refuse}

\label{lesson:HI-C349-inkarkarna}



\begin{quote}
Before the new one: say the Hindi for an accusation, then the Hindi for to blurt out.
\end{quote}

\subsection*{You'll want to know: \hi{इनकार} \hi{करना}}
\textbf{\hi{इनकार} \hi{करना}} — \emph{inkār karnā} — \textquotedblleft{}to deny, to refuse\textquotedblright{}.

\begin{quote}
\hi{नेता} \hi{ने} \hi{हर} \hi{सवाल} \hi{का} \hi{जवाब} \hi{देने} \hi{से} \hi{इनकार} \hi{किया।}
\end{quote}

\begin{quote}
\emph{netā ne har savāl kā javāb dene se inkār kiyā.}
\end{quote}

\begin{quote}
The politician refused to answer any question.
\end{quote}

Past, with \hi{ने} (he or she …): \hi{उसने} \hi{इनकार} \hi{किया}, \emph{usne inkār kiyā}.

\begin{cousinweb}
\textbf{\hi{इनकार} \hi{करना}} (\emph{inkār karnā}) joins \hi{इनकार} to \hi{करना}, \textquotedblleft{}to do\textquotedblright{}. \hi{इनकार} is usually traced to Arabic \emph{inkār}, \textquotedblleft{}denial\textquotedblright{}, which reached Hindi through Persian.
\end{cousinweb}

\begin{grammarlens}[title={What you've built}]
One more word for giving reasons and opinions.
\end{grammarlens}

\subsection*{Your turn}
\begin{itemize}
\raggedright
  \item \emph{Say it:} \emph{inkār karnā}
  \item \emph{Say it:} \emph{inkār karnā}, once more
  \item \emph{Say it:} \emph{netā ne har savāl kā javāb dene se inkār kiyā.} — The politician refused to answer any question.
  \item \emph{Recall:} say the Hindi for an accusation, then the Hindi for to blurt out, then say \emph{inkār karnā} again
\end{itemize}

\begin{tcolorbox}[breakable,colback=teal!4,colframe=teal!35!black,title={Before you move on}]
What does \hi{इनकार} \hi{करना} mean? (\textquotedblleft{}To deny, to refuse\textquotedblright{}.) Say it once more.
\end{tcolorbox}
\section[\texorpdfstring{\hi{पुष्टि} \hi{करना} (puṣṭi karnā)}{पुष्टि करना (puṣṭi karnā)}]{\hi{पुष्टि} \hi{करना} (puṣṭi karnā) — to confirm}

\label{lesson:HI-C349-pustikarna}



\begin{quote}
Before the new one: say the Hindi for to blurt out, then the Hindi for to deny.
\end{quote}

\subsection*{You'll want to know: \hi{पुष्टि} \hi{करना}}
\textbf{\hi{पुष्टि} \hi{करना}} — \emph{puṣṭi karnā} — \textquotedblleft{}to confirm\textquotedblright{}.

\begin{quote}
\hi{डाक्टर} \hi{ने} \hi{ख़बर} \hi{की} \hi{पुष्टि} \hi{की।}
\end{quote}

\begin{quote}
\emph{ḍākṭar ne khabar kī puṣṭi kī.}
\end{quote}

\begin{quote}
The doctor confirmed the news.
\end{quote}

Past, with \hi{ने} (he or she …): \hi{उसने} \hi{ख़बर} \hi{की} \hi{पुष्टि} \hi{की}, \emph{usne khabar kī puṣṭi kī}.

\begin{cousinweb}
\textbf{\hi{पुष्टि} \hi{करना}} (\emph{puṣṭi karnā}) joins \hi{पुष्टि} to \hi{करना}, \textquotedblleft{}to do\textquotedblright{}. \hi{पुष्टि} is a tatsama, borrowed from Sanskrit \emph{puṣṭi}, \textquotedblleft{}growth, thriving\textquotedblright{} (Monier-Williams). The Hindi sense, a confirmation, is not among the senses Monier-Williams gives.
\end{cousinweb}

\begin{grammarlens}[title={What you've built}]
One more word for giving reasons and opinions.
\end{grammarlens}

\subsection*{Your turn}
\begin{itemize}
\raggedright
  \item \emph{Say it:} \emph{puṣṭi karnā}
  \item \emph{Say it:} \emph{puṣṭi karnā}, once more
  \item \emph{Say it:} \emph{ḍākṭar ne khabar kī puṣṭi kī.} — The doctor confirmed the news.
  \item \emph{Recall:} say the Hindi for to blurt out, then the Hindi for to deny, then say \emph{puṣṭi karnā} again
\end{itemize}

\begin{tcolorbox}[breakable,colback=teal!4,colframe=teal!35!black,title={Before you move on}]
What does \hi{पुष्टि} \hi{करना} mean? (\textquotedblleft{}To confirm\textquotedblright{}.) Say it once more.
\end{tcolorbox}
\section[\texorpdfstring{\hi{स्वीकार} \hi{करना} (svīkār karnā)}{स्वीकार करना (svīkār karnā)}]{\hi{स्वीकार} \hi{करना} (svīkār karnā) — to admit, to own up to}

\label{lesson:HI-C349-svikarkarna}



\begin{quote}
Before the new one: say the Hindi for to deny, then the Hindi for to confirm.
\end{quote}

\subsection*{You'll want to know: \hi{स्वीकार} \hi{करना}}
\textbf{\hi{स्वीकार} \hi{करना}} — \emph{svīkār karnā} — \textquotedblleft{}to admit, to own up to\textquotedblright{}.

\begin{quote}
\hi{अरुण} \hi{ने} \hi{अपनी} \hi{ग़लती} \hi{स्वीकार} \hi{की।}
\end{quote}

\begin{quote}
\emph{Aruṇ ne apnī ġaltī svīkār kī.}
\end{quote}

\begin{quote}
Arun admitted his mistake.
\end{quote}

Past, with \hi{ने} (he or she …): \hi{उसने} \hi{ग़लती} \hi{स्वीकार} \hi{की}, \emph{usne ġaltī svīkār kī}.

\begin{cousinweb}
\textbf{\hi{स्वीकार} \hi{करना}} (\emph{svīkār karnā}) joins \hi{स्वीकार} to \hi{करना}, \textquotedblleft{}to do\textquotedblright{}. \hi{स्वीकार} is a tatsama, borrowed from Sanskrit \emph{svīkāra}, \textquotedblleft{}making one's own, acceptance\textquotedblright{} (Monier-Williams).
\end{cousinweb}

\begin{grammarlens}[title={What you've built}]
One more word for giving reasons and opinions.
\end{grammarlens}

\subsection*{Your turn}
\begin{itemize}
\raggedright
  \item \emph{Say it:} \emph{svīkār karnā}
  \item \emph{Say it:} \emph{svīkār karnā}, once more
  \item \emph{Say it:} \emph{Aruṇ ne apnī ġaltī svīkār kī.} — Arun admitted his mistake.
  \item \emph{Recall:} say the Hindi for to deny, then the Hindi for to confirm, then say \emph{svīkār karnā} again
\end{itemize}

\begin{tcolorbox}[breakable,colback=teal!4,colframe=teal!35!black,title={Before you move on}]
What does \hi{स्वीकार} \hi{करना} mean? (\textquotedblleft{}To admit, to own up to\textquotedblright{}.) Say it once more.
\end{tcolorbox}
\section[\texorpdfstring{\hi{ग़लतफ़हमी} (ġalatfahmī)}{ग़लतफ़हमी (ġalatfahmī)}]{\hi{ग़लतफ़हमी} (ġalatfahmī) — a misunderstanding}

\label{lesson:HI-C349-galatfahmi}



\begin{quote}
Before the new one: say the Hindi for to confirm, then the Hindi for to admit.
\end{quote}

\subsection*{You'll want to know: \hi{ग़लतफ़हमी}}
\textbf{\hi{ग़लतफ़हमी}} — \emph{ġalatfahmī} — \textquotedblleft{}a misunderstanding\textquotedblright{}.

\begin{quote}
\hi{यह} \hi{बहस} \hi{एक} \hi{ग़लतफ़हमी} \hi{की} \hi{वजह} \hi{से} \hi{हुई।}
\end{quote}

\begin{quote}
\emph{yah bahas ek ġalatfahmī kī vajah se huī.}
\end{quote}

\begin{quote}
This argument happened because of a misunderstanding.
\end{quote}

Feminine.

\begin{cousinweb}
\textbf{\hi{ग़लतफ़हमी}} (\emph{ġalatfahmī}) is usually traced to Persian \emph{ġalat-fahmī}, \textquotedblleft{}mistaken understanding\textquotedblright{}. Both parts are Arabic words that Persian put together.
\end{cousinweb}

\begin{grammarlens}[title={What you've built}]
One more word for giving reasons and opinions.
\end{grammarlens}

\subsection*{Your turn}
\begin{itemize}
\raggedright
  \item \emph{Say it:} \emph{ġalatfahmī}
  \item \emph{Say it:} \emph{ġalatfahmī}, once more
  \item \emph{Say it:} \emph{yah bahas ek ġalatfahmī kī vajah se huī.} — This argument happened because of a misunderstanding.
  \item \emph{Recall:} say the Hindi for to confirm, then the Hindi for to admit, then say \emph{ġalatfahmī} again
\end{itemize}

\begin{tcolorbox}[breakable,colback=teal!4,colframe=teal!35!black,title={Before you move on}]
What does \hi{ग़लतफ़हमी} mean? (\textquotedblleft{}A misunderstanding\textquotedblright{}.) Say it once more.
\end{tcolorbox}
